\section{Motivation And Workloads}
\label{sec:workload-characterization}

A VFS name-resolution extension targets a narrow workload pattern in which the
source system changes which existing object a pathname should select, but the
operation does not need a new filesystem object model.
The source-system evidence below identifies that pattern and separates it from
the broader behavior that should remain with agent runtimes, environment
builders, services, FUSE filesystems, or custom filesystems.

\subsection{State-Dependent Path Views}

Under this definition, a transition qualifies when a system state change alters
which existing object a pathname denotes, or whether the object is visible,
while ordinary lower-filesystem semantics remain unchanged. The definition is
intentionally narrower than "programmable filesystem." It
excludes synthetic file contents, custom metadata persistence, distributed
directory indexing, data-path mediation, write-conflict resolution, and
cross-path transactional snapshots.

Within this definition, the transition is the workload state change, the policy
is the eBPF decision logic, and the action is the bounded result that the kernel
validates for lookup or directory enumeration.

The distinction matters because the natural implementation boundary changes.
If the system needs synthetic contents or custom metadata persistence, a FUSE
service, custom filesystem, metadata service, or application/runtime mechanism
is the appropriate implementation boundary. If the state change only affects
lookup or directory visibility, a VFS name-resolution policy can express the
relevant policy. The remaining source-system behavior stays outside name
resolution.

\subsection{Workload Sources}

We use source systems to choose workloads and correctness conditions, not as
separate contributions. We call each extracted correctness condition a source
oracle. A same-oracle KVM run applies that condition unchanged to \namei and its
comparison rows.

The first source family is agent/workspace state. AgentFS is the primary source
because its public implementation exposes workspace lifecycle traces
from which path-selection and visibility oracles can be extracted~\cite{agentfs_repo}.
SWE-Factory-Gym supplies a released Click source task whose fail-to-pass test
can run through those workspace views~\cite{swe_factory_repo,swe_factory_gym_dataset}.
BranchFS, Sandlock, and YoloFS provide neighboring branch, sandbox, staging,
and stackable-filesystem context rather than separate evaluation
workloads~\cite{branchfs_repo,sandlock_repo,yolofs_repo}.

The traditional source families are application file sharing, build action
sandboxing, checkpoint/restore, HPC staging, service rotation, and toolchain
environments. The application-sharing workload follows the XDG Documents
portal grant/revoke behavior~\cite{xdg_documents_portal}. The build-action
workload follows Bazel's action-specific input views~\cite{bazel_sandboxing}.
The toolchain workload follows per-environment software selection and
switch/rollback behavior represented by profile systems such as
Nix~\cite{nix_store}. Kubernetes AtomicWriter supplies the completed
already-materialized payload-view sequence~\cite{kubernetes_projected_volumes}.
DMTCP path virtualization supplies an additional source-defined lifecycle
shape but remains a portfolio case until its exact KVM oracle
runs~\cite{dmtcp_path_virtualization}. Full projected-volume materialization
and service reload also remain outside the completed Kubernetes subset.

A historical ccache matrix supplies supporting compile-time evidence, but is
not a headline source workload: ccache already implements cache lookup and
validation in userspace. Source systems are admitted only when the behavior
under study is the application's path view, not an application-level cache
algorithm re-expressed in BPF.

Only source behavior tied to a concrete same-oracle KVM run becomes final
evaluation evidence. Source inspection and paper-derived behavior motivate
workload shape and boundary comparisons.

\begin{table}[t]
\centering
\scriptsize
\begin{tabular}{@{}L{0.2\linewidth}L{0.25\linewidth}L{0.25\linewidth}L{0.2\linewidth}@{}}
\hline
Example & State and operation & \namei action and delegated behavior & Oracle \\
\hline
Agent workspace lifecycle & branch or checkpoint state changes lookup and
readdir under a workspace path & select or hide existing lower objects while
the source runtime keeps COW, checkpoint, sandbox, writes, and lifecycle
orchestration & final tree, lookup/readdir coherence, and preserved
lower-filesystem permissions \\
Application file sharing & application identity and grant state change whether
one existing host document is visible & select the granted document or hide it
while grant persistence and portal UI remain outside the hook & grant, revoke,
cross-application isolation, and unchanged lower object \\
Build action sandboxing & action identity and declared-input set choose which
existing input tree one logical pathname exposes & select the action's declared
object or hide an undeclared name while Bazel keeps process and build semantics
& real action success, distinct concurrent outputs, and undeclared-input
failure \\
\hline
\end{tabular}
\caption{Representative state-dependent path-view examples. The \namei
evidence is the lookup/readdir action, and non-name-resolution behavior remains
with the source runtime or lower filesystem.}
\label{tab:path-view-examples}
\end{table}

Three design constraints follow from the source systems. First, each workload
carries its own state, such as a branch identifier, checkpoint, cache epoch,
build generation, service epoch, or secret epoch. Policy therefore must run at
the affected lookup or directory-enumeration operation and remain scoped per
workload.

Second, the path-view portion of each source system is narrower than the full
system. Many rows also include runtime orchestration, service reload,
synthetic contents, or storage semantics that remain separate responsibilities.
The narrower path-view scope implies that the kernel and lower filesystem must
keep data, permission, and persistence ownership.

Third, directory visibility and lookup selection appear together in the source
oracles. Lookup and readdir therefore need one coherent decision contract.

Together, these source-system observations set the paper's comparison rule:
(1) source behavior supplies the oracle, (2) feature-equivalent FUSE supplies
the direct programmable-policy cost comparison, and (3) custom or stackable
filesystems define the broader filesystem-method and runtime-responsibility
boundary. Materialized
namespace mechanisms remain related-work and background comparisons rather than
the central experiment. The same rule drives
the design requirements in Section~\ref{sec:design} and the RQ-organized
evaluation in Section~\ref{sec:evaluation}.

\subsection{Selected Workloads}

The evaluated RQ1 set contains five non-overlapping workflows. The two main
cases are Agent Workspaces and Build Action Sandboxing. The formal Agent
workspace case exercises branch-visible staged/base selection, hides,
readdir coherence, and lower-tree mutations. It also runs a released Click
source task through concurrent base/completed views, switch, rollback, and
withdrawal. The formal build-action run executes two concurrent real Bazel
actions with distinct declared-input roots and an undeclared-input failure.
The formal application-sharing run executes official
\code{xdg-document-portal} 1.18.4 and \namei under the same
per-application grant/isolation/revoke oracle. The formal Kubernetes case follows the official
\code{AtomicWriter} payload-view sequence through update, no-op, and rollback
under one stable logical root. The formal toolchain case selects two existing
CPython environments, switches one process-group view, and rolls it back.
These three supporting workflows broaden the source-system evidence without
becoming separate headline claims.

Full service validation and reload, checkpoint/restore, and HPC file staging
remain portfolio extensions rather than completed experiments. A workload
becomes evidence only after a reviewed KVM same-oracle run. In every completed
case, policy runs at lookup and directory enumeration; non-name-resolution
behavior stays with the source system or lower filesystem.
